\section{HyNoC Architecture}
\label{sec:architecture}

\subsection{Overview}

A HyNoC network is built by assembling routers with a configurable number of ports
(3 to 9 in the current implementation). Each router port is full-duplex, composed
of an ingress interface and an egress interface. Two neighboring routers are
connected by cross-linking their egress and ingress interfaces. Processing nodes
are attached through a \emph{local interface}, which adds an extra output FIFO to
decouple the node's clock domain from the router.

The router implements a \textbf{full crossbar}: any ingress port can simultaneously
communicate with any egress port, subject to arbitration. Multiple paths can be
open concurrently inside the same router, enabling full bisection bandwidth
utilization when traffic permits. Figure~\ref{fig:noc-overview} shows the internal
organization of 5-port routers and their interconnections.

\begin{figure}[t]
  \centering
  \includegraphics[width=\columnwidth]{hynoc_noc_overview}
  \caption{Router interconnections overview.}
  \label{fig:noc-overview}
\end{figure}

\subsection{Ingress Port}

The ingress port receives flits from the egress of a neighboring router. It buffers
incoming flits in a dual-clock FIFO (configurable depth: 2 to 64 entries), which
simultaneously handles clock domain crossing between the upstream router's clock
domain and the local router domain.

Once flits are available, the ingress FSM decodes the routing flit, extracts the
next-hop egress port identifier, and asserts a \emph{request} signal to the
corresponding egress port. It then waits for a \emph{grant}. Upon grant, the
ingress begins forwarding flits: if the routing flit still carries hops for
downstream routers (index $> 0$), it is updated (index decremented) and forwarded;
if the index has reached zero, the routing flit is consumed locally and not
forwarded. Subsequent payload flits are streamed directly to the granted egress
until the stop bit of the last flit is detected, at which point the channel is
released.

Flow control is managed by monitoring the FIFO level of the downstream router's
ingress FIFO. The ingress stalls when the downstream FIFO approaches full,
preventing flit loss. The internal data path of a 3-port router is illustrated in
Figure~\ref{fig:data-arch}.

\begin{figure}[t]
  \centering
  \includegraphics[width=\columnwidth]{hynoc_router_data_arch}
  \caption{Internal data path of a 3-port router.}
  \label{fig:data-arch}
\end{figure}

\subsection{Egress Port}

The egress port receives concurrent write requests from all $N-1$ ingress ports.
It arbitrates among them using the PRRA, grants one request at a time, and
connects the selected ingress data path to the downstream router's ingress FIFO.
The downstream FIFO level is forwarded back to the granted ingress for flow
control. The control path is shown in Figure~\ref{fig:ctrl-arch}: contrary to the
data path, control signals are point-to-point between each ingress and each egress,
avoiding broadcast of control information.

\begin{figure}[t]
  \centering
  \includegraphics[width=\columnwidth]{hynoc_router_control_arch}
  \caption{Internal control path of a 3-port router.}
  \label{fig:ctrl-arch}
\end{figure}

\subsection{Parallel Round-Robin Arbiter (PRRA)}

The standard sequential round-robin FSM introduces up to $N$ cycles of latency
before granting a request. HyNoC replaces this with a \textbf{parallel
round-robin arbiter}: a LUT-based FSM that evaluates all pending requests
simultaneously, selects the highest-priority pending request according to a
rotating priority order, and transitions directly to the corresponding state.

The state transition $P_k$ is given by:
\begin{equation}
  P_k \leftarrow \text{request}[k] = 1
  \label{eq:prra}
\end{equation}
where $P_0$ carries the highest priority and $P_{N-2}$ the lowest. An additional
implicit transition keeps the current state while a granted request is active,
preventing spurious re-arbitration during packet transfer.

The arbiter responds in \textbf{one or two clock cycles} (configurable via the
\texttt{PRRA\_PIPELINE} parameter), regardless of the number of ports. The
critical path does not grow with $N$ because each state's transition equation is
stored in a dedicated LUT, selected by a multiplexer from the state register.
Starvation is prevented by the rotating priority: after serving request $k$,
priority rotates to port $k+1$.

\subsection{Clock Domain Architecture}

Each router port operates in its own clock domain. The dual-clock FIFO at each
ingress interface handles the crossing between the upstream port clock and the
local router clock. For timing-critical designs, a single-clock mode
(\texttt{SINGLE\_CLOCK\_ROUTER=1}) is available, reducing crossing latency.

The local interface adds an extra FIFO at the egress output, fully decoupling the
attached node's clock domain from the router.

\subsection{First Layer Protocol}

\subsubsection{Packet and Flit Structure}

The network operates on flits of $(\mathtt{PAYLOAD\_WIDTH}+1)$ bits. The most
significant bit is a \emph{stop bit}: when set, it marks the last flit of a
packet. A packet consists of one or more \emph{routing flits} followed by one or
more \emph{payload flits}. The routing flit carries a 4-bit \emph{Proto} field:

\begin{table}[h]
\centering
\caption{Supported routing protocols}
\label{tab:proto}
\begin{tabular}{@{}ll@{}}
\toprule
Proto & Routing method \\
\midrule
\texttt{4'b0000} & Unicast circuit switch \\
\texttt{4'b0001} & Multicast circuit switch \\
\texttt{4'b1000} & XY routing (reserved) \\
\texttt{4'b1111} & Forbidden \\
\bottomrule
\end{tabular}
\end{table}

\subsubsection{Unicast Source Routing}

The routing flit carries $H$ hop fields of $\lceil\log_2(P-1)\rceil$ bits each,
encoding the relative egress port ID at the next router. An \emph{index} field
(initialized to $H-1$) points to the hop to be consumed by the next router. The
ingress port decrements the index after consuming a hop; if the index was already
zero, the routing flit is not forwarded downstream. The relative hop encoding is
illustrated in Figure~\ref{fig:hops}.

\begin{figure}[t]
  \centering
  \includegraphics[width=0.85\columnwidth]{hynoc_hops_encoding}
  \caption{Relative hop encoding for a $3\times 3$ mesh network.}
  \label{fig:hops}
\end{figure}

The relative encoding (egress IDs numbered counterclockwise from the ingress port,
excluding that port) reduces both the hop field width and the internal crossbar
size.

\subsubsection{Multicast Source Routing}

In multicast mode, each hop field is a bitmask of width $P-1$, where each bit
selects one egress port. The packet is duplicated and forwarded to all asserted
egress ports simultaneously. Multiple routing flits with different multicast masks
can be chained to address complex multicast trees. This is particularly useful for
broadcast operations in VLIW many-core workloads.
